\clearpage % On s'assure d'être sur la bonne page
\phantomsection % Crée une ancre invisible pour le lien du sommaire
\section*{Introduction} % Le titre dans le document (sans numéro)
\addcontentsline{toc}{section}{Introduction} % Ajout manuel au sommaire
\markboth{Introduction}{Introduction} % Pour que le titre apparaisse en haut des pages suivantes


Dans l'écosystème financier contemporain, la vitesse n'est plus un simple avantage compétitif mais la condition essentielle de l'existence des acteurs de marché. Cette course effrénée à la vitesse a propulsé le trading haute fréquence au devant de la scène. Les décisions de vente et d'achat se prennent alors en quelques nanosecondes. Un élément crucial des systèmes financiers modernes est le \textbf{\textit{Order Matching Engine}} qui fait correspondre des ordres d'achat à des ordres de vente pour exécuter une transaction. Une priorité selon le prix et le temps est alors effectuée. L'ordre d'achat le plus élevé pour un ordre de vente donné est priorisé, et à prix égal, l'ordre ayant été effectué en premier a la priorité. On voit alors de manière évidente les enjeux liés à la vitesse d'exécution.

Dans ce cadre, le C++ s'est imposé comme langage de programmation de référence pour ses possibilités d'optimisation bas niveau (gestion manuelle de la mémoire et intéraction directe avec le matériel). À l'inverse, Java a mis du temps à trouver sa place au sein de l'univers de l'ultra faible latence. Son modèle de gestion automatique de la mémoire pilotée par le garbage collector crée des pauses imprévisibles qui augmentent fortement la latence dans les pires cas, une situation inacceptable pour les marchés financiers.

Une nouvelle philosophie propulsée par Martin Thompson, co-fondateur de LMAX, la mechanical sympathy propose que les performances d'un programme dépendent de la capacité du programmeur à le développer en s'alignant sur les contraintes matérielles du CPU et de la mémoire. Je vais dans ce mémoire explorer comment, en appliquant les concepts de mechanical sympathy, on peut transformer la JVM en environnement déterministe et à faible latence pour tirer partie du riche écosystème de Java tout en rivalisant avec les performances du C++ au niveau de la latence. La problématique sera donc: \textbf{Comment l’alignement de l’architecture logicielle sur les contraintes matérielles et le tuning de l’environnement d’exécution permettent-ils de transformer la Java Virtual Machine en une plateforme déterministe et distribuée pour le trading haute fréquence ?}

Pour répondre à cette problématique, mon plan s'articulera en deux parties:

\begin{itemize}
  \item Un état de l'art: j'analyserai la hiérarchie des caches CPU, les obstacles internes de la JVM (Just In Time, Safepoints, Stop-the-World, \dots) et nous étudierons les solutions industrielles telles que LMAX Disruptor, Aeron et Chronicle Software
  \item Une implémentation incrémentale d'un \textit{order matching engine}: je détaillerai la construction de ma solution à travers dix jalons qui montrent l'évolution de la performance des différentes versions du matching engine en partant d'une base qui s'appuie sur des principes de programmation orientée objet puis nous appliqueront une architecture Zéro Garbage Collector et Data Oriented ainsi que des optimisations systèmes qui permettront de montrer les gains apportés par l'application de la mechanical sympathy à notre problématique.
\end{itemize}

Mon intérêt pour cette problématique prend racine dans une passion profonde pour l’ingénierie logicielle. Je suis fermement convaincu que l’excellence d’un ingénieur réside dans sa maîtrise de la \textit{mechanical sympathy}. Au-delà de cet attrait personnel, ce mémoire s'inscrit dans un enjeu industriel majeur. Dans des environnements de finance de marché où les capitaux exposés au risque se chiffrent en millions d’euros à chaque microseconde, l’optimisation n’est plus une option, mais une nécessité économique vitale. Ce projet est pour moi l'opportunité de démontrer que la rigueur de l'implémentation est un levier majeur de création de valeur financière.